\begin{recipe}
[% 
    preparationtime = {\unit[1]{h}\unit[30]{min} + \unit[24]{h} Kühlzeit + \unit[24]{h} Marinierzeit},
    bakingtime = {\unit[75]{min} Dämpfen + \unit[15]{min} Anbraten},
    portion = {\portion{6}},
    source = {\protect\recipesource{https://www.zuckerjagdwurst.com/de/rezepte/veganer-sauerbraten-mit-spaetzle-und-rotkraut}{Zucker\&Jagdwurst: Veganer Sauerbraten}},
    calory = {405kcal | 64g Protein | 13g KH | 11g Fett},
    price = {2,35€},
    created = {30.05.2026},
    rating = {4},
]
{Seitan}

%\graph{% pictures
%    big=pic/basics/07_seitan-sauerbraten
%}

\introduction{%
    Man kann Seitan auch aus Mehl auswaschen, aber es macht einen das Leben einfacher den Seitan-Fix zu nehmen. Seitan wird oft als Ersatzprodukt für Ente benutzt.
}

\ingredients{
    \textbf{Seitanbraten} & \\
    500 g & Seitan-Fix\index{Tofu \& Fleischalternativen!Seitan}\\
    500 ml & Kalte Gemüsebrühe\\
    2 EL & Sojasauce\\
    2 EL & Senf\\
    2 EL & Hefeflocken\\
    2 TL & Paprikapulver\\
    2 TL & Knoblauchpulver\\
    1 TL & Zwiebelpulver\\
    1 TL & Thymian\\
    1 TL & Pfeffer\\
    1 TL & Salz\\
    1 Prise & Muskat\\
    1 Prise & Curry
}

\preparation{%
    \step%
    Seitan-Fix mit Paprikapulver, Knoblauchpulver, Zwiebelpulver, Thymian, Pfeffer, Salz, Hefeflocken, Muskat und Curry in einer großen Schüssel vermischen.
    
    \step%
    Kalte Gemüsebrühe, Sojasauce und Senf zugeben und alles 5--10 Minuten kräftig durchkneten, bis eine kompakte, elastische Masse entsteht.
    
    \step%
    Die Masse zu einem großen Braten formen und 30 Minuten bei Zimmertemperatur ruhen lassen.
    
    \step%
    In Gemüsebrühe den Seitan ca 20 Minuten auf niedriger Stufe kochen lassen, das Wasser sollte nicht sprudeln, sondern nur leicht simmern.
    
    \step%
    Pflanzenöl in einer großen Pfanne oder einem Bräter erhitzen und den Seitan von allen Seiten kräftig anbraten, bis er schön gebräunt ist.
}

\hint{%
    Für die Weiterverwendung den angebratenen Seitan in Scheiben schneiden oder mit zwei Gabeln grob zerreißen und anschließend in einer Bratensauce oder in Mushroomsauce ziehen lassen. Gerade das Zerreißen hilft, damit mehr Sauce in den Seitan kommt.
}

\end{recipe}